\documentclass[11pt,a4paper]{article}

% ── Packages ──────────────────────────────────────────────────────────────────
\usepackage[utf8]{inputenc}
\usepackage[T1]{fontenc}
\usepackage{amsmath,amssymb,amsthm}
\usepackage{booktabs}
\usepackage{graphicx}
\usepackage{geometry}
\usepackage{hyperref}
\usepackage{xcolor}
\usepackage{algorithm}
\usepackage{algpseudocode}
\usepackage{caption}
\usepackage{subcaption}
\usepackage{multirow}
\usepackage{array}

\geometry{margin=1in}
\hypersetup{colorlinks=true, linkcolor=blue!60!black, citecolor=green!50!black, urlcolor=blue!70!black}

\newcommand{\E}{\mathbb{E}}
\newcommand{\R}{\mathbb{R}}
\newcommand{\N}{\mathbb{N}}
\DeclareMathOperator*{\argmax}{arg\,max}

\title{\textbf{Phase 3: Computational Findings Report} \\[6pt]
\large Dynamic Augmented Pooled Testing Strategies (DAPTS)}
\author{Héctor Becerril Villamil \and Francisco Marmolejo-Cossío \and Nicholas Lopez}
\date{March 2026}

\begin{document}
\maketitle

\begin{abstract}
This document consolidates the computational findings from Phase~3 of the DAPTS project.
We report results from three sprints:
(1)~Gibbs sampler verification and a critical mixing-failure bug fix,
(2)~the $\beta$-meta\-parameter for balancing myopic clearing reward with information gain, and
(3)~large-scale experiments ($n=20$--$50$) comparing greedy strategies.
Key findings include a fundamental mixing pathology in Gibbs sampling under overlapping exact-count constraints,
a prevalence-dependent regime where $\beta>0$ improves expected utility by up to 26\%, and
empirical gap analyses for MOSEK-based greedy at scale.
\end{abstract}

\tableofcontents
\newpage

%% ═══════════════════════════════════════════════════════════════════════════════
\section{Background: The Augmented Model}
%% ═══════════════════════════════════════════════════════════════════════════════

We consider a population of $n$ individuals, each independently infected with prior probability $p_i$.
A \emph{pooled test} on subset $t \subseteq [n]$ returns the \textbf{exact count} of infected:
\[
  r = |t \cap Z|, \qquad Z = \{i : X_i = 1\}.
\]
The decision-maker has $B$ tests, each of maximum pool size $G$.
Clearing individual $i$ (certifying $X_i=0$) yields utility $u_i$.
The objective is to maximize $\E\bigl[\sum_{i \in \text{cleared}} u_i\bigr]$ over adaptive testing strategies.

The \textbf{Bayesian posterior} after test history $H_k = ((t_1,r_1),\ldots,(t_k,r_k))$ is computed by:
\begin{itemize}
  \item \textbf{Exact counting} (brute-force over $2^n$ infection profiles): $O(2^n)$, feasible for $n \leq 14$.
  \item \textbf{Gibbs sampling}: MCMC approximation scaling to $n \sim 50{+}$.
\end{itemize}


%% ═══════════════════════════════════════════════════════════════════════════════
\section{Sprint 1: Gibbs Sampler Verification}
\label{sec:gibbs}
%% ═══════════════════════════════════════════════════════════════════════════════

\subsection{Systematic Comparison Protocol}

We compared Gibbs posterior marginals against exact counting across 6 configurations $\times$ 20 random instances = 120 test cases:

\begin{table}[h]
\centering
\caption{Test configurations for Gibbs vs.\ exact comparison.}
\label{tab:gibbs_configs}
\begin{tabular}{cccc}
\toprule
\textbf{Config} & $n$ & $B$ & $G$ \\
\midrule
1 & 5 & 2 & 3 \\
2 & 5 & 3 & 3 \\
3 & 6 & 2 & 3 \\
4 & 6 & 3 & 4 \\
5 & 7 & 2 & 4 \\
6 & 7 & 3 & 4 \\
\bottomrule
\end{tabular}
\end{table}

For each instance: prior $p_i \sim \text{Uniform}(0.05, 0.4)$, utility $u_i \sim \text{Uniform}(1, 10)$, a random infection profile $Z$ sampled from the prior, 1--2 greedy tests executed, and the posterior compared. Pass criterion: $\max_i |q^{\text{exact}}_i - q^{\text{Gibbs}}_i| < 0.05$.


\subsection{The Mixing Failure Bug}

\begin{finding}
\textbf{Before fix:} fraction passed $= 0.858$, worst-case error $= 0.799$.
\end{finding}

\paragraph{Root cause.}
When multiple tests produce overlapping exact-count constraints, the feasible state space can decompose into \textbf{disconnected components}.
Consider two pools $t_1 = \{0,1,2\}$ and $t_2 = \{2,3,4\}$ with results $r_1 = 1, r_2 = 1$.
The constraint that $|t_1 \cap Z| = 1$ \emph{and} $|t_2 \cap Z| = 1$ allows states like $Z = \{0,3\}$ and $Z = \{1,4\}$, but a single-site Gibbs flip cannot move from one to the other without passing through an infeasible intermediate state.

\paragraph{Standard Gibbs iteration.}
At each step, the sampler picks agent $i$ and samples $X_i \mid X_{-i}$ from its conditional:
\[
  P(X_i = 1 \mid X_{-i}) = \frac{p_i \cdot \mathbf{1}[\text{state valid with } X_i=1]}
  {p_i \cdot \mathbf{1}[\text{valid, }X_i{=}1] + (1-p_i) \cdot \mathbf{1}[\text{valid, }X_i{=}0]}.
\]
If flipping any single agent violates a constraint, the chain is \emph{stuck}: it has zero probability of ever reaching the other component.


\subsection{The Fix: Exact Fallback + Global Block Moves}

We implemented a two-pronged fix in \texttt{bayesian.py}:

\begin{enumerate}
  \item \textbf{Exact fallback} for $|\text{active set}| \leq 7$:
    Since $2^7 = 128$ profiles is trivial to enumerate, we bypass Gibbs entirely and compute the exact posterior by counting.
    This covers all test configurations in Table~\ref{tab:gibbs_configs}.

  \item \textbf{Global pairwise Metropolis--Hastings block moves} for $|\text{active set}| > 7$:
    After each Gibbs sweep, propose $\max(n_{\text{active}}, 5)$ swaps of the form $(i_{\text{inf}} \leftrightarrow j_{\text{healthy}})$
    drawn uniformly from the \emph{full} active set (not restricted to a single pool).
    Accept with MH ratio:
    \[
      \alpha = \min\!\left(1,\; \frac{p_j (1-p_i)}{p_i (1-p_j)}\right)
    \]
    provided the proposed state satisfies all constraints.
    This allows the chain to jump between disconnected components when the swap preserves all test counts.
\end{enumerate}

\begin{algorithm}[h]
\caption{Gibbs update with exact fallback and block moves}
\label{alg:gibbs_fix}
\begin{algorithmic}[1]
\Require Prior $p$, history $H_k$, population size $n$
\State Deterministic deductions (clear pools with $r=0$, confirm all-infected pools)
\State Compute active set $A$ from remaining undetermined constraints
\If{$|A| \leq 7$}
    \State \Return \textsc{ExactCounting}$(p, H_k, n)$ \Comment{$O(2^7)$ trivial}
\EndIf
\State Initialize state by sampling from priors
\For{iteration $= 1, \ldots, T$}
    \For{each agent $i \in A$ (random order)} \Comment{Standard Gibbs}
        \State Sample $X_i \mid X_{-i}$ from conditional
    \EndFor
    \For{each test pool} \Comment{Per-pool swap moves}
        \State Propose infected$\leftrightarrow$healthy swap within pool; MH accept/reject
    \EndFor
    \For{$m = 1, \ldots, \max(|A|, 5)$} \Comment{Global block moves}
        \State Pick random infected $i$, healthy $j$ from $A$
        \State Swap $X_i \leftrightarrow X_j$; validate all constraints; MH accept/reject
    \EndFor
    \If{iteration $>$ burn-in}
        \State Accumulate sample counts
    \EndIf
\EndFor
\State \Return Empirical marginals from accumulated samples
\end{algorithmic}
\end{algorithm}


\subsection{Results After Fix}

\begin{table}[h]
\centering
\caption{Gibbs vs.\ exact comparison: before and after fix.}
\label{tab:gibbs_before_after}
\begin{tabular}{lcc}
\toprule
\textbf{Metric} & \textbf{Before} & \textbf{After} \\
\midrule
Fraction passed ($\epsilon < 0.05$) & 0.858 & \textbf{1.000} \\
Worst-case max error & 0.799 & \textbf{0.000} \\
Instances tested & 120 & 120 \\
\bottomrule
\end{tabular}
\end{table}

The exact fallback completely eliminates approximation error for $n_\text{active} \leq 7$, which covers all configurations in the test suite.
For larger active sets, the global block moves improve mixing but may not fully resolve all pathological configurations---an open direction for future work (see Section~\ref{sec:discussion}).

\subsection{Expected Utility Comparison}

A secondary test compared Gibbs-based greedy expected utility against counting-based expected utility (4 configs $\times$ 10 instances = 40 tests).
Pass criterion: relative error $< 5\%$.
\textbf{Result: 40/40 passed} after the fix.


%% ═══════════════════════════════════════════════════════════════════════════════
\section{Sprint 2: The $\beta$ Meta-Parameter}
\label{sec:beta}
%% ═══════════════════════════════════════════════════════════════════════════════

\subsection{Formulation}

The standard myopic greedy selects the pool maximizing immediate expected clearing reward:
\[
  t^* = \argmax_{t \subseteq A,\; |t| \leq G} \; P(r{=}0 \mid t) \cdot \sum_{i \in t} u_i.
\]
We introduce a meta-parameter $\beta \geq 0$ that adds an information-gain bonus:
\begin{equation}
\label{eq:beta_score}
  \text{Score}(t) = P(r{=}0 \mid t) \cdot \sum_{i \in t} u_i \;+\; \beta \cdot \E[\text{InfoGain}(t)],
\end{equation}
where InfoGain can be measured as entropy reduction, variance reduction, or count of newly confirmed agents.

\paragraph{Intuition.}
When $\beta = 0$, the strategy is purely myopic---it only values pools likely to clear immediately.
When $\beta > 0$, the strategy trades some immediate reward for tests that resolve uncertainty about the population, potentially enabling better decisions in future rounds.


\subsection{VIP Benchmark: Two Prevalence Regimes}

We test on a structured VIP scenario ($n=20$: 8 VIP agents with $u_i=10$, 12 regular with $u_i=2$):

\begin{table}[h]
\centering
\caption{VIP-high-prevalence scenario: $p_{\text{VIP}}=0.8$, $p_{\text{reg}}=0.2$, $B=6$, $G=10$.}
\label{tab:vip_high}
\begin{tabular}{rrrr}
\toprule
$\beta$ & $\E[U]$ & First pool & Pool size \\
\midrule
0.0 &  8.50 & $\{0\}$ & 1 \\
0.1 & 10.00 & $\{0\}$ & 1 \\
0.5 & 10.00 & $\{0\}$ & 1 \\
1.0 & 10.00 & $\{0\}$ & 1 \\
2.0 & 14.00 & $\{0\}$ & 1 \\
5.0 &  8.00 & $\{0,1,2,3\}$ & 4 \\
10.0 &  6.00 & $\{0,1,2,3\}$ & 4 \\
\bottomrule
\end{tabular}
\end{table}

\begin{finding}
At high prevalence ($p=0.8$), $P(r{=}0)$ drops \textbf{exponentially} with pool size:
$P(r{=}0) = \prod_{i \in t}(1-p_i) = 0.2^{|t|}$.
For $|t|=4$: $P(r{=}0) = 0.0016$.
The myopic clearing term dominates $\beta \cdot \text{InfoGain}$ for moderate $\beta$, so the greedy always selects individual tests ($|t|=1$).
Only at $\beta \geq 5$ does information gain override, but this \emph{hurts} expected utility.
\end{finding}

\begin{table}[h]
\centering
\caption{VIP-moderate-prevalence scenario: $p_{\text{VIP}}=0.35$, $p_{\text{reg}}=0.1$, $B=6$, $G=5$.}
\label{tab:vip_moderate}
\begin{tabular}{rrrr}
\toprule
$\beta$ & $\E[U]$ & First pool & Pool size \\
\midrule
0.0 & 39.50 & $\{0,1\}$ & 2 \\
0.1 & 40.00 & $\{0,1\}$ & 2 \\
0.5 & 40.00 & $\{0,1\}$ & 2 \\
1.0 & 33.00 & $\{0,8,9,10,11\}$ & 5 \\
2.0 & \textbf{41.50} & $\{0,8,9,10,11\}$ & 5 \\
5.0 & 32.00 & $\{8,9,10,11,12\}$ & 5 \\
10.0 & 29.50 & $\{8,9,10,11,12\}$ & 5 \\
\bottomrule
\end{tabular}
\end{table}

\begin{finding}
At moderate prevalence ($p=0.35$), the information-gain bonus at $\beta=2.0$ shifts the first pool from a size-2 VIP-only test to a size-5 mixed pool including regular agents.
This yields $\E[U] = 41.50$, a \textbf{5\% improvement} over $\beta=0$ ($\E[U]=39.50$).
The mixed pool leverages the low-prevalence regular agents ($p=0.1$) to increase $P(r{=}0)$ while still covering a VIP agent.
\end{finding}


\subsection{Randomized Beta Sweep}

Across 2 randomized VIP instances (per-instance seeds) with 5 $\beta$ values:

\begin{table}[h]
\centering
\caption{Beta sweep: mean $\pm$ std of expected utility over randomized instances.}
\label{tab:beta_sweep}
\begin{tabular}{llcc}
\toprule
\textbf{Scenario} & $\beta$ & Mean $\E[U]$ & Std \\
\midrule
\multirow{5}{*}{VIP ($p_{\text{VIP}}\approx0.8$)}
 & 0.0 & 14.43 & 0.72 \\
 & 0.5 & 15.90 & 1.75 \\
 & 1.0 & \textbf{18.19} & 1.05 \\
 & 2.0 & 17.77 & 0.32 \\
 & 5.0 & 18.11 & 0.28 \\
\midrule
\multirow{5}{*}{Uniform ($p_i \in [0.1,0.4]$)}
 & 0.0 & 6.38 & 1.88 \\
 & 0.5 & 6.38 & 1.88 \\
 & 1.0 & 6.38 & 1.88 \\
 & 2.0 & 6.38 & 1.88 \\
 & 5.0 & 5.63 & 1.88 \\
\bottomrule
\end{tabular}
\end{table}

\begin{finding}
For the VIP scenario, $\beta = 1.0$ achieves a \textbf{26\% improvement} over $\beta=0$ ($18.19$ vs.\ $14.43$).
For the uniform scenario, $\beta$ has no effect up to $\beta=2.0$ and is harmful at $\beta=5.0$.
This suggests the $\beta$ meta-parameter is most valuable when the population has \textbf{heterogeneous utilities and moderate prevalence}.
\end{finding}


%% ═══════════════════════════════════════════════════════════════════════════════
\section{Sprint 3: Large-Scale Experiments}
\label{sec:sprint3}
%% ═══════════════════════════════════════════════════════════════════════════════

\subsection{Main Configuration Results}

Quick-run results with MOSEK-optimized greedy pool selection ($n=20, B=2, G=10$):

\begin{table}[h]
\centering
\caption{Config D ($n=20$, $B=2$, $G=10$, medium prevalence): MOSEK greedy vs.\ baselines.}
\label{tab:sprint3_main}
\begin{tabular}{lrl}
\toprule
\textbf{Strategy} & $\E[U]$ & \textbf{Note} \\
\midrule
$U_{\max}$ (clairvoyant upper bound) & 103.71 & Theoretical maximum \\
$U_{\text{greedy,MOSEK}}$ & 36.09 & MOSEK exponential-cone solver \\
$U_{\text{greedy,enum}}$ & 36.09 & Exhaustive enumeration (validates MOSEK) \\
$U_{\beta\text{-greedy}}$ ($\beta=1$) & 30.11 & Beta-greedy with info gain \\
$U_{\text{single}}$ & 15.64 & Individual testing only \\
\bottomrule
\end{tabular}
\end{table}

\begin{finding}
MOSEK and exhaustive enumeration produce \textbf{identical results} at $n=20$, validating the solver.
Greedy pooling captures $\frac{36.09 - 15.64}{103.71 - 15.64} = 23.2\%$ of the gap between single testing and the clairvoyant optimum.
With only $B=2$ tests, the gap to $U_{\max}$ remains large---this motivates investigating larger $B$.
\end{finding}


\subsection{VIP Experiment Results}

\begin{table}[h]
\centering
\caption{VIP config V1: $n=20$ (8 VIP + 12 regular), $B=6$, $G=10$.}
\label{tab:sprint3_vip}
\begin{tabular}{lrl}
\toprule
\textbf{Strategy} & $\E[U]$ & \textbf{Note} \\
\midrule
$U_{\max}$ & 72.92 & \\
$U_{\text{greedy,MOSEK}}$ & 46.09 & \\
$U_{\beta\text{-greedy}}$ ($\beta=1$) & 45.00 & \\
$U_{\text{single}}$ & 39.34 & \\
\bottomrule
\end{tabular}
\end{table}

Greedy pooling captures $\frac{46.09 - 39.34}{72.92 - 39.34} = 20.1\%$ of the gap.
The $\beta$-greedy is slightly below standard greedy in this single instance; a larger sample is needed for robust comparison.


\subsection{Utility Distribution Modulation}

\begin{table}[h]
\centering
\caption{Effect of utility distribution ($n=20$, $B=5$, $G=10$).}
\label{tab:utility_mod}
\begin{tabular}{lccccc}
\toprule
\textbf{Distribution} & $U_{\max}$ & $U_{\text{greedy}}$ & $U_{\beta}$ & $U_{\text{single}}$ & Greedy gap \\
\midrule
Uniform ($u_i=1$) & 15.55 & 7.83 & 5.25 & 4.20 & 31.9\% \\
Skewed ($u_i \sim \text{Pareto}$) & 93.19 & 58.04 & 61.25 & 41.46 & 32.1\% \\
\bottomrule
\end{tabular}
\end{table}

\begin{finding}
With skewed utilities, $\beta$-greedy ($\E[U]=61.25$) \textbf{outperforms} standard greedy ($\E[U]=58.04$) by 5.5\%.
This confirms that when a few individuals have disproportionately high utility, the information-gain bonus helps the strategy identify and prioritize them.
Under uniform utilities, $\beta$-greedy underperforms---additional exploration is wasteful when all individuals are equally valuable.
\end{finding}


\subsection{Large Pool Size ($G$) Exploration}

Francisco's question: ``With only $B=2$ tests, what happens as we increase $G$?''

\begin{table}[h]
\centering
\caption{$\E[U]$ vs.\ maximum pool size $G$ ($n=20$, $B=2$, medium prevalence).}
\label{tab:large_G}
\begin{tabular}{rccc}
\toprule
$G$ & $U_{\text{greedy,MOSEK}}$ & $U_{\max}$ & Gap \\
\midrule
5  & 27.60 & 97.03 & 71.6\% \\
10 & 27.60 & 97.03 & 71.6\% \\
\bottomrule
\end{tabular}
\end{table}

\begin{finding}
Increasing $G$ from 5 to 10 does not improve $\E[U]$ for $B=2$ in this instance.
The myopic greedy already selects a pool of optimal size $\leq 5$; allowing larger pools does not help because $P(r{=}0)$ decreases with pool size faster than the utility sum increases.
This suggests that the \textbf{effective pool size is self-limiting} under the myopic strategy.
\end{finding}


%% ═══════════════════════════════════════════════════════════════════════════════
\section{Overnight Baseline Results ($n=10$)}
\label{sec:overnight}
%% ═══════════════════════════════════════════════════════════════════════════════

Prior to the Phase~3 sprints, we ran baseline experiments at $n=10$ with both Gurobi MILP and exhaustive enumeration:

\begin{table}[h]
\centering
\caption{Overnight baseline: $n=10$, $B=2$, $G=3$, Gurobi solver.}
\label{tab:overnight}
\begin{tabular}{llccccc}
\toprule
$n$ & Regime & $U_{\max}$ & $U_{\text{single}}$ & $U_{\text{greedy}}$ & Time (s) & Gap \\
\midrule
10 & low  & 41.80 & 14.21 & 31.56 & 0.073 & 24.5\% \\
10 & low  & 54.69 & 16.26 & 37.97 & 0.005 & 30.6\% \\
10 & med  & 35.78 & 12.32 & 20.71 & 0.005 & 35.6\% \\
10 & med  & 46.77 & 13.72 & 24.26 & 0.006 & 31.8\% \\
10 & high & 25.12 &  8.99 & 10.16 & 0.005 & 59.6\% \\
10 & high & 32.76 &  9.50 & 10.40 & 0.005 & 68.3\% \\
\midrule
20 & low  & 120.49 & 18.37 & 46.41 & 0.004 & 61.5\% \\
20 & low  & 109.83 & 17.62 & 45.87 & 0.005 & 58.2\% \\
20 & med  & 103.71 & 15.64 & 31.37 & 0.005 & 69.8\% \\
20 & med  &  93.07 & 15.23 & 29.50 & 0.004 & 68.3\% \\
20 & high &  73.55 & 11.61 & 14.28 & 0.006 & 80.6\% \\
20 & high &  64.10 & 11.04 & 13.13 & 0.007 & 79.5\% \\
\bottomrule
\end{tabular}
\end{table}

\begin{finding}
The gap between greedy and $U_{\max}$ \textbf{grows sharply with prevalence}: from $\sim$25--30\% (low) to $\sim$60--80\% (high).
At high prevalence, $P(r{=}0)$ is small for any non-trivial pool, so greedy barely outperforms individual testing.
This reinforces the motivation for the $\beta$ meta-parameter: at high prevalence, pure myopic reward is a poor objective.
\end{finding}


%% ═══════════════════════════════════════════════════════════════════════════════
\section{MOSEK Solver Validation}
\label{sec:mosek}
%% ═══════════════════════════════════════════════════════════════════════════════

The MOSEK exponential-cone formulation for myopic pool selection solves:
\[
  \max_{x \in \{0,1\}^n} \sum_i u_i x_i + \sum_i \ln(1-p_i) \cdot x_i
  \quad \text{s.t.} \quad \sum_i x_i \leq G, \quad x_i = 0 \;\forall i \notin A
\]
via the relaxation with exponential cone constraints for $\ln(1-p_i)$.

\begin{itemize}
  \item \textbf{Correctness}: MOSEK and exhaustive enumeration produce identical pool selections at $n=20$ (Table~\ref{tab:sprint3_main}).
  \item \textbf{Performance}: $\sim$6.5s per pool selection at $n=20$, $B=5$, $G=10$. Enumeration takes $\sim$11.6s (Table~\ref{tab:sprint3_main}).
  \item \textbf{Scalability}: Tested up to $n=50$, $G=10$ (Sprint~3 Config~C).
\end{itemize}


%% ═══════════════════════════════════════════════════════════════════════════════
\section{Discussion and Open Questions}
\label{sec:discussion}
%% ═══════════════════════════════════════════════════════════════════════════════

\subsection{Gibbs Mixing for Large Active Sets}

The exact fallback for $|A| \leq 7$ resolves the mixing failure for all tested configurations.
However, for $|A| > 7$ with heavily overlapping constraints, the global pairwise block moves may not fully connect all feasible components.
We observed persistent errors (max error $= 0.93$) in a synthetic test with $n=10$, two overlapping pools of size 5.

\textbf{Potential improvements:}
\begin{itemize}
  \item \emph{Multi-site block moves}: swap $k > 2$ agents simultaneously.
  \item \emph{Parallel tempering}: run chains at different ``temperatures'' and exchange states.
  \item \emph{Exact sub-block sampling}: decompose the active set into blocks coupled by shared pool membership and sample blocks exactly when small enough.
\end{itemize}

For practical use in the greedy algorithm, the posterior update after each individual test uses the exact single-test formula (\texttt{bayesian\_update\_single\_test}), which is $O(n)$ and exact.
Gibbs is only needed when computing the joint posterior over multiple overlapping tests, which occurs in the brute-force DP solver---not in the myopic greedy.

\subsection{$\beta$ Regime Characterization}

Our findings suggest a \textbf{two-regime characterization}:
\begin{itemize}
  \item \textbf{Low/moderate prevalence} ($p \lesssim 0.4$): $P(r{=}0)$ is non-negligible for pools of size 3--5. The $\beta$ bonus can steer the strategy toward larger, more informative pools. Optimal $\beta \in [1, 2]$.
  \item \textbf{High prevalence} ($p \gtrsim 0.6$): $P(r{=}0)$ collapses exponentially. Individual tests dominate regardless of $\beta$. Only extreme $\beta$ ($\geq 5$) overrides this, but at a cost to expected utility.
\end{itemize}

\subsection{Scaling Bottleneck: The DP Solver}

The brute-force DP solver (for computing $U_{\max}$) becomes intractable beyond $n \approx 10$:
the state space involves frozensets of $2^n$ infection profiles, causing $>90$\,GB memory at $n=10$.
This limits our ability to compute optimality gaps at scale.
Future work should investigate information-theoretic bounds or LP relaxations for $U_{\max}$.

\subsection{Next Steps (Sprint 4: Paper Examples)}

\begin{enumerate}
  \item Hand-computable examples for the Overleaf paper ($n=3,4$ with exact DP trees).
  \item Run Sprint~3 experiments at full scale ($\texttt{--n-instances 50}$) for statistically robust gap estimates.
  \item Investigate adaptive $\beta$ schedules: start with high $\beta$ (exploration) and decay to $0$ (exploitation) across the $B$ rounds.
\end{enumerate}


%% ═══════════════════════════════════════════════════════════════════════════════
\section{Summary of Key Findings}
\label{sec:summary}
%% ═══════════════════════════════════════════════════════════════════════════════

\begin{enumerate}
  \item \textbf{Gibbs mixing failure} (Sprint~1): Overlapping exact-count constraints create disconnected feasible regions. Fixed with exact fallback ($|A| \leq 7$) and global MH block moves. All 120 test cases now pass.

  \item \textbf{$\beta$-meta-parameter} (Sprint~2): At moderate prevalence, $\beta \in [1,2]$ improves $\E[U]$ by up to 26\% for VIP scenarios. At high prevalence, $\beta$ is ineffective or harmful. Under skewed utilities, $\beta$-greedy outperforms standard greedy by 5.5\%.

  \item \textbf{Large-scale greedy} (Sprint~3): MOSEK solver validated against enumeration at $n=20$. Greedy captures 20--32\% of the gap to $U_{\max}$ depending on utility distribution. The gap grows with prevalence (25\% low $\to$ 80\% high).

  \item \textbf{Self-limiting pool size}: Increasing $G$ beyond $\sim$5 yields no improvement for the myopic greedy at moderate prevalence, because $P(r{=}0)$ drops faster than utility increases.

  \item \textbf{Prevalence is the key regime parameter}: All findings---greedy gap, $\beta$ effectiveness, pool size selection---are governed by the interaction between prevalence and $P(r{=}0)$.
\end{enumerate}


\end{document}
